% EM structure and teaching revision 20260928
\subsection{相位谐振、幅值峰值与参数依赖}
串联理想RLC接恒幅电压源时，$|I|=U/\sqrt{R^2+(\omega L-1/(\omega C))^2}$，电抗相消使电流最大且输入电压、电流同相。此处阻尼只来自与频率无关的串联电阻，所以两种谐振判据重合。
含有损耗线圈的并联电路则未必如此。线圈为 $r+\ii\omega L$，与电容并联时
\[
 Y=\frac{r}{r^2+\omega^2L^2}
 +\ii\left(\omega C-\frac{\omega L}{r^2+\omega^2L^2}\right).
\]
令虚部为零可求相位谐振频率，但实部本身也随频率变化，因此 $|Y|$ 的极小值不必出现在同一点。若问电源电流最小，应直接极小化 $|Y|$；若问电压电流同相，才使用 $\Im Y=0$。
比较谐振曲线还必须固定激励方式。恒压源下看电流响应，恒流源下看电压响应，峰谷可能互换。品质因数很高时一些频率差别很小，可以作近似，但应先给出近似参数，例如 $r\ll\omega L$，再说明忽略的是哪一阶修正。

